\chapter{Literature Review}

\section{Self-Supervised Speech Representations (Wav2Vec2)}
The landscape of Automatic Speech Recognition (ASR) has been revolutionized by self-supervised learning, drastically reducing the reliance on massive, human-transcribed datasets. At the forefront of this shift is the Wav2Vec2 framework \cite{baevski2020wav2vec}. Wav2Vec2 operates by encoding raw audio waveforms into latent speech representations and subsequently utilizing a Transformer network to build contextualized representations across the entire utterance. During its pre-training phase, it solves a contrastive task, learning to identify the true quantized latent representation for a masked time step among a set of distractors. 

This self-supervised pre-training yields a robust acoustic feature extractor capable of capturing deep phonetic information without requiring explicit labels. For the VoiceScore project, Wav2Vec2 serves as the foundational backbone, allowing the model to leverage vast amounts of unlabeled speech to understand general acoustic properties before being fine-tuned on the specific, narrow phonetic inventory of Indian English.

\section{Connectionist Temporal Classification (CTC)}
Traditionally, training ASR models required frame-level alignment between the audio input and the target phonemes, a labor-intensive and error-prone prerequisite. The introduction of Connectionist Temporal Classification (CTC) \cite{graves2006connectionist} eliminated this requirement. CTC introduces a 'blank' token and computes the probability of a target sequence by marginalizing over all possible valid frame-level alignments. 

While highly effective for general ASR, standard CTC models are prone to producing "peaky" probability distributions, where the model outputs high confidence for a phoneme over a very narrow time window and predicts the blank token elsewhere. This peakiness complicates fine-grained pronunciation scoring, as it relies on stable frame-level probabilities across the duration of a spoken sound.

\section{Goodness of Pronunciation (GoP)}
Goodness of Pronunciation (GoP) is the de facto standard metric for Computer-Assisted Pronunciation Training (CAPT) systems. Originally formalized by Witt and Young \cite{witt2000gop}, GoP is defined as the log-likelihood ratio between the probability of the target phoneme and the most likely alternative phoneme, given the acoustic segment. 

In classical systems based on Gaussian Mixture Models and Hidden Markov Models (GMM-HMM), calculating GoP required explicit pre-segmentation of the audio to define the boundaries of each phoneme. However, with the advent of deep learning and CTC, recent research has focused on adapting GoP calculations to operate directly on the posterior probabilities generated by neural networks without requiring strict pre-segmentation. By extracting the log-probabilities (posteriors) for each phoneme at each time step, it is possible to derive GoP scores that accurately reflect the speaker's acoustic precision relative to the expected canonical sequence.

\section{Contrastive Learning in ASR}
Contrastive learning has gained traction as a mechanism to improve the robustness and confidence calibration of ASR models \cite{drexler2023contrastive}. In the context of CTC, a Contrastive CTC objective can be employed to explicitly manage the embedding space of the acoustic features and the target vocabulary. 

By adding a contrastive loss alongside the standard CTC training, the model is encouraged to project similar speech segments (e.g., various instances of a retroflex stop) closer together in a shared metric space, while pushing dissimilar phonemes apart. This acts as a powerful regularizer, making the encoder invariant to nuisance factors such as background noise or minor regional accent variations, while remaining highly sensitive to the critical phonetic distinctions required for accurate pronunciation diagnosis. This contrastive projection mechanism forms a core pillar of the VoiceScore architecture.

\section{Phonological Characteristics of Indian English}
Indian English represents a distinct, institutionalized variety of English with its own well-defined phonological rules \cite{pandey2021indian}. The phonetic inventory diverges significantly from standard Received Pronunciation (RP) or General American (GA). Key characteristics include the absence of alveolar stops (/t/, /d/), which are universally replaced by retroflex stops (/ʈ/, /ɖ/). Dental fricatives (/θ/, /ð/) are typically realized as dental plosives (/t̪/, /d̪/). Furthermore, the distinction between the labiodental fricative /v/ and the bilabial approximant /w/ is frequently neutralized into a single labiodental approximant /ʋ/. Understanding and formalizing these phonological rules was critical for developing the tailored lexicon required to train an unbiased diagnostic ASR.
